\documentclass[../../数学分析培优讲义(答案版).tex]{subfiles}

\begin{document}

\setcounter{chapter}{3}
\pagestyle{fancy}

\chapter{一元函数微分学}

\section{导数}

\subsection{导数的定义及可微性验证}

\begin{example}
设 $f(x)=|x|$,则 $f(x)$ 在 $x\neq0$ 处可导,且
\[
    f'(x)=\dfrac{|x|}{x}.
\]
\end{example}

\begin{proof}[解]
当 $x>0$ 时 $f(x)=x$，当 $x<0$ 时 $f(x)=-x$。因此对 $x\neq0$，
\[
    f'(x)=
    \begin{cases}
        1, & x>0,\\
        -1, & x<0,
    \end{cases}
    =\dfrac{|x|}{x}.
\]
\end{proof}

\begin{example}
证明:$f(x)=|x|^3$ 在 $x=0$ 处的三阶导数 $f'''(0)$ 不存在.
\end{example}

\begin{proof}[解]
分别在原点两侧求导，可得
\[
    f(x)=
    \begin{cases}
        x^3, & x\geq0,\\
        -x^3, & x<0,
    \end{cases}
    \qquad
    f''(x)=6|x|.
\]
特别地，$f''(0)=0$。但
\[
    \lim\limits_{h\to0^+}\dfrac{f''(h)-f''(0)}{h}=6,
    \qquad
    \lim\limits_{h\to0^-}\dfrac{f''(h)-f''(0)}{h}=-6.
\]
左右极限不相等，故 $f'''(0)$ 不存在。
\end{proof}

\begin{example}
设
\[
f(x)=
\begin{cases}
    x^2\sin\dfrac{\pi}{x}, & x<0,\\
    A, & x=0,\\
    ax^2+b, & x>0,
\end{cases}
\]
其中 $A,a,b$ 为常数,试问 $A,a,b$ 为何值时 $f(x)$ 在 $x=0$ 处可导,为什么?并求 $f'(0)$.
\end{example}

\begin{proof}[解]
可导必连续。由左右极限
\[
    \lim\limits_{x\to0^-}x^2\sin\dfrac{\pi}{x}=0,
    \qquad
    \lim\limits_{x\to0^+}(ax^2+b)=b,
\]
可知连续性要求且只要求 $A=b=0$。此时
\[
    \lim\limits_{x\to0^-}\dfrac{f(x)-f(0)}{x}
    =\lim\limits_{x\to0^-}x\sin\dfrac{\pi}{x}=0,
\]
而
\[
    \lim\limits_{x\to0^+}\dfrac{f(x)-f(0)}{x}
    =\lim\limits_{x\to0^+}ax=0.
\]
因而 $a$ 可取任意实数，$A=b=0$，且 $f'(0)=0$。
\end{proof}

\begin{example}
设函数 $f(x)$ 在闭区间 $\lbrack 0,1\rbrack$ 上四次连续可微,$f(0)=0,f'(0)=0$,证明函数
\[
F(x)=
\begin{cases}
    \dfrac{f(x)}{x^2}, & 0<x\leq1,\\
    \dfrac{f''(0)}{2}, & x=0,
\end{cases}
\]
在 $\lbrack 0,1\rbrack$ 上二次连续可微.
\end{example}

\begin{proof}[解]
由 $f(0)=f'(0)=0$ 及 Taylor 公式的积分余项，对 $x\in\lbrack0,1\rbrack$ 有
\[
    f(x)=x^2\int_0^1(1-t)f''(tx)\,\mathrm{d}t.
\]
因此可将 $F$ 统一写成
\[
    F(x)=\int_0^1(1-t)f''(tx)\,\mathrm{d}t.
\]
当 $x=0$ 时右端等于 $f''(0)/2$，与题设定义一致。由于 $f\in C^4(\lbrack0,1\rbrack)$，可在积分号下连续求导两次：
\[
    F'(x)=\int_0^1t(1-t)f'''(tx)\,\mathrm{d}t,
\]
\[
    F''(x)=\int_0^1t^2(1-t)f^{(4)}(tx)\,\mathrm{d}t.
\]
上述三个积分都是 $x$ 的连续函数，故 $F\in C^2(\lbrack0,1\rbrack)$。
\end{proof}

\begin{example}
设 $f(x)$ 在 $x_0$ 的某邻域内有定义.
\begin{enumerate}[label=(\arabic*)]
    \item 若 $f(x)$ 在 $x_0$ 处可导,试证
    \[
        \lim\limits_{h\to0}
        \dfrac{f(x_0+h)-f(x_0-h)}{2h}=f'(x_0);
    \]
    \item 反之,若上式左端极限存在,是否能推出 $f'(x_0)$ 存在?若结论成立,请证明;不成立请给出反例.
\end{enumerate}
\end{example}

\begin{proof}[解]
\begin{enumerate}[label=(\arabic*)]
    \item 将差商拆成
    \[
        \dfrac{f(x_0+h)-f(x_0-h)}{2h}
        =\dfrac{1}{2}\dfrac{f(x_0+h)-f(x_0)}{h}
        +\dfrac{1}{2}\dfrac{f(x_0)-f(x_0-h)}{h}.
    \]
    当 $h\to0$ 时，两项分别趋于 $f'(x_0)/2$，故原极限为 $f'(x_0)$。
    \item 不能。取 $f(x)=|x-x_0|$，则对任意 $h\neq0$，
    \[
        \dfrac{f(x_0+h)-f(x_0-h)}{2h}=0,
    \]
    故对称差商的极限存在；但 $f$ 在 $x_0$ 处的左右导数分别为 $-1$ 和 $1$，因而不可导。
\end{enumerate}
\end{proof}

\begin{example}[\markstar]
证明 Riemann 函数
\[
R(x)=
\begin{cases}
    \dfrac{1}{q}, & x=\dfrac{p}{q},\\
    0, & x=0,1\text{ 或无理数},
\end{cases}
\]
处处不可微.
\end{example}

\begin{proof}[解]
先设 $x_0=p_0/q_0\in(0,1)$ 为既约有理数。取无理数列 $x_n\to x_0$，则
\[
    \left|\dfrac{R(x_n)-R(x_0)}{x_n-x_0}\right|
    =\dfrac{1}{q_0|x_n-x_0|}\longrightarrow+\infty,
\]
故 $R$ 在 $x_0$ 处不可导。

再设 $x_0\in(0,1)$ 为无理数。由 Dirichlet 逼近定理，存在无穷多个既约分数 $p/q$，使
\[
    0<\left|x_0-\dfrac{p}{q}\right|<\dfrac{1}{q^2}.
\]
沿这些有理点趋于 $x_0$ 时，
\[
    \left|\dfrac{R(p/q)-R(x_0)}{p/q-x_0}\right|
    =\dfrac{1/q}{|p/q-x_0|}>q,
\]
差商不可能收敛。

最后，在端点 $0$ 处分别沿无理数和 $x=1/q$ 趋近，右差商分别为 $0$ 和 $1$；在端点 $1$ 处分别沿无理数和 $x=(q-1)/q$ 趋近，左差商分别为 $0$ 和 $-1$。因而 $R$ 处处不可微。
\end{proof}

\begin{example}
求 $f(x)=[x]\sin\pi x$ 的单侧导数,并讨论可微性.
\end{example}

\begin{proof}[解]
当 $x\in(m,m+1)$，$m\in\mathbb{Z}$ 时，$[x]=m$，故
\[
    f'(x)=m\pi\cos\pi x.
\]
对整数 $n$，$f(n)=0$。当 $h\to0^+$ 时 $[n+h]=n$，当 $h\to0^-$ 时 $[n+h]=n-1$，因而
\[
    f'_+(n)=n\pi(-1)^n,
    \qquad
    f'_-(n)=(n-1)\pi(-1)^n.
\]
二者不相等，故 $f$ 在每个整数点都不可微，在其余各点可微。
\end{proof}

\subsection{高阶导数}

\methodsection{(A) 利用已知结论}

\begin{example}
计算 $y^{(n)}$,设
\begin{enumerate}[label=(\arabic*)]
    \item $y=\sin^6x+\cos^6x$;
    \item $y=\sin ax\cdot\sin bx$;
    \item $y=\dfrac{x^2+x+1}{x^2-5x+6}$.
\end{enumerate}
\end{example}

\begin{proof}[解]
\begin{enumerate}[label=(\arabic*)]
    \item 由
    \[
        \sin^6x+\cos^6x
        =1-3\sin^2x\cos^2x
        =\dfrac{5}{8}+\dfrac{3}{8}\cos4x,
    \]
    当 $n\geq1$ 时
    \[
        y^{(n)}=\dfrac{3}{8}4^n
        \cos\left(4x+\dfrac{n\pi}{2}\right).
    \]
    \item 积化和差得
    \[
        y=\dfrac{1}{2}\bigl[\cos(a-b)x-\cos(a+b)x\bigr],
    \]
    因而
    \[
        y^{(n)}=\dfrac{1}{2}\left[
        (a-b)^n\cos\left((a-b)x+\dfrac{n\pi}{2}\right)
        -(a+b)^n\cos\left((a+b)x+\dfrac{n\pi}{2}\right)
        \right].
    \]
    \item 部分分式分解为
    \[
        y=1-\dfrac{7}{x-2}+\dfrac{13}{x-3}.
    \]
    故当 $n\geq1$ 时
    \[
        y^{(n)}=(-1)^n n!\left[
        -\dfrac{7}{(x-2)^{n+1}}+\dfrac{13}{(x-3)^{n+1}}
        \right].
    \]
\end{enumerate}
\end{proof}

\methodsection{(B) 利用 Leibniz 公式}

\begin{example}
设 $f_1(x),f_2(x)$ 在 $x_0$ 及其附近有定义,在 $x_0$ 处有直到 $n$ 阶导数,记
\[
    N(f)=\sum\limits_{k=0}^n\dfrac{1}{k!}|f^{(k)}(x_0)|.
\]
试证明:
\[
    N(f_1f_2)\leq N(f_1)N(f_2).
\]
\end{example}

\begin{proof}[解]
由 Leibniz 公式，对 $0\leq j\leq n$ 有
\[
    |(f_1f_2)^{(j)}(x_0)|
    \leq\sum\limits_{k=0}^{j}C_{k}^{j}
    |f_1^{(k)}(x_0)|\,|f_2^{(j-k)}(x_0)|.
\]
因此
\[
\begin{aligned}
    N(f_1f_2)
    &\leq\sum\limits_{j=0}^{n}\sum\limits_{k=0}^{j}
    \dfrac{|f_1^{(k)}(x_0)|}{k!}
    \dfrac{|f_2^{(j-k)}(x_0)|}{(j-k)!}\\
    &=\sum\limits_{\substack{k,l\geq0\\k+l\leq n}}
    \dfrac{|f_1^{(k)}(x_0)|}{k!}
    \dfrac{|f_2^{(l)}(x_0)|}{l!}\\
    &\leq N(f_1)N(f_2).
\end{aligned}
\]
\end{proof}

\begin{example}
试证:
\[
    (a^2+b^2)^{\frac{n}{2}}e^{ax}\sin(bx+n\varphi)
    =\sum\limits_{i=0}^n C_{n}^{i}a^{n-i}b^i e^{ax}
    \sin\left(bx+\dfrac{\pi i}{2}\right),
\]
其中
\[
    \varphi=\arctan\dfrac{b}{a}.
\]
\end{example}

\begin{proof}[解]
令 $r=\sqrt{a^2+b^2}$，并选取角 $\varphi$ 使
\[
    \cos\varphi=\dfrac{a}{r},
    \qquad
    \sin\varphi=\dfrac{b}{r}.
\]
直接归纳可得
\[
    \dfrac{\mathrm{d}^n}{\mathrm{d}x^n}
    \bigl(e^{ax}\sin bx\bigr)
    =r^n e^{ax}\sin(bx+n\varphi).
\]
另一方面，对 $e^{ax}\sin bx$ 使用 Leibniz 公式，得
\[
\begin{aligned}
    \dfrac{\mathrm{d}^n}{\mathrm{d}x^n}
    \bigl(e^{ax}\sin bx\bigr)
    &=\sum\limits_{i=0}^{n}C_{i}^{n}
    a^{n-i}e^{ax}b^i
    \sin\left(bx+\dfrac{\pi i}{2}\right).
\end{aligned}
\]
比较两式，并代入 $r^n=(a^2+b^2)^{n/2}$，即得所证恒等式。
\end{proof}

\methodsection{(C) 利用数学归纳法}

\begin{example}[\markstar]
证明:若函数
\[
f(x)=
\begin{cases}
    e^{-\frac{1}{x^2}}, & x\neq0,\\
    0, & x=0,
\end{cases}
\]
则在 $x=0$ 处,
\[
    f^{(n)}(0)=0,
    \qquad n=1,2,\ldots.
\]
\end{example}

\begin{proof}[解]
当 $x\neq0$ 时，对任意 $n\geq0$，反复求导可将 $f^{(n)}$ 写成
\[
    f^{(n)}(x)=P_n\left(\dfrac{1}{x}\right)e^{-1/x^2},
\]
其中 $P_n$ 是多项式。这由归纳立即得到：若对 $n$ 成立，则求导后仍是 $1/x$ 的多项式乘以 $e^{-1/x^2}$。

对任意正整数 $m$，都有
\[
    \lim\limits_{x\to0}|x|^{-m}e^{-1/x^2}=0.
\]
因而若已知 $f^{(n)}(0)=0$，则
\[
    f^{(n+1)}(0)
    =\lim\limits_{h\to0}\dfrac{f^{(n)}(h)-f^{(n)}(0)}{h}
    =\lim\limits_{h\to0}\dfrac{P_n(1/h)e^{-1/h^2}}{h}=0.
\]
由 $f(0)=0$ 归纳得 $f^{(n)}(0)=0$，$n=1,2,\ldots$。
\end{proof}

\begin{example}[推广的 Leibniz]
设 $u_1(x),u_2(x),\ldots,u_k(x)$ 都是 $x$ 的 $n$ 次可微函数,试证:
\[
    (u_1u_2\cdots u_k)^{(n)}
    =\sum\limits_{r_1+\cdots+r_k=n}
    \dfrac{n!}{r_1!r_2!\cdots r_k!}
    u_1^{(r_1)}\cdots u_k^{(r_k)}.
\]
\end{example}

\begin{proof}[解]
对 $k$ 作归纳。$k=2$ 时即为 Leibniz 公式。假设结论对 $k-1$ 个因子成立，令
$v=u_1u_2\cdots u_{k-1}$，则
\[
\begin{aligned}
    (vu_k)^{(n)}
    &=\sum\limits_{s=0}^{n}C_{s}^{n}v^{(s)}u_k^{(n-s)}\\
    &=\sum\limits_{s=0}^{n}C_{s}^{n}
      \sum\limits_{r_1+\cdots+r_{k-1}=s}
      \dfrac{s!}{r_1!\cdots r_{k-1}!}
      u_1^{(r_1)}\cdots u_{k-1}^{(r_{k-1})}u_k^{(n-s)}.
\end{aligned}
\]
令 $r_k=n-s$，并利用
\[
    C_s^n\dfrac{s!}{r_1!\cdots r_{k-1}!}
    =\dfrac{n!}{r_1!\cdots r_k!},
\]
即得题设公式。
\end{proof}

\begin{example}[\markstar]
$f(x)$ 有任意阶导数,$f'(x)=(f(x))^2$,求 $f^{(n)}(x)\ (n>2)$.
\end{example}

\begin{proof}[解]
对 $n$ 作归纳。当 $n=1$ 时，题设已给出
\[
    f'(x)=f(x)^2=1!f(x)^2.
\]
若 $f^{(n)}(x)=n!f(x)^{n+1}$，则
\[
\begin{aligned}
    f^{(n+1)}(x)
    &=n!(n+1)f(x)^n f'(x)\\
    &=(n+1)!f(x)^{n+2}.
\end{aligned}
\]
因此对所有 $n\geq1$，
\[
    f^{(n)}(x)=n!f(x)^{n+1}.
\]
\end{proof}

\methodsection{(D) 利用递推公式}

\begin{example}
设 $f(x)=(\arcsin x)^2$,求 $f^{(n)}(0)$.
\end{example}

\begin{proof}[解]
令 $y=(\arcsin x)^2$。直接求导可得微分方程
\[
    (1-x^2)y''-xy'=2.
\]
记 $a_n=y^{(n)}(0)$。对上式比较 $x^m$ 的系数，或对方程求 $m$ 次导后令 $x=0$，得
\[
    a_{m+2}=m^2a_m,
    \qquad m\geq1.
\]
又 $a_0=a_1=0$，$a_2=2$。因而所有奇数阶导数都为 $0$，而对 $k\geq1$，
\[
    a_{2k}=2\cdot2^2\cdot4^2\cdots(2k-2)^2
    =2^{2k-1}[(k-1)!]^2.
\]
即
\[
    f^{(n)}(0)=
    \begin{cases}
        0, & n\text{ 为奇数},\\
        2^{n-1}\left[\left(\dfrac{n}{2}-1\right)!\right]^2,
        & n\text{ 为正偶数}.
    \end{cases}
\]
\end{proof}

\begin{example}
证明 Legendre 多项式
\[
    p_n(x)=\dfrac{1}{2^n n!}\bigl[(x^2-1)^n\bigr]^{(n)},
    \qquad n=0,1,2,\ldots,
\]
满足方程
\[
    (1-x^2)p_n''(x)-2xp_n'(x)+n(n+1)p_n(x)=0.
\]
\end{example}

\begin{proof}[解]
令 $u(x)=(x^2-1)^n$，则
\[
    (x^2-1)u'(x)=2nx\,u(x).
\]
对上式两端同时求 $n+1$ 次导。由于 $x^2-1$ 与 $x$ 的高阶导数很快为零，Leibniz 公式给出
\[
    (x^2-1)u^{(n+2)}+2x u^{(n+1)}-n(n+1)u^{(n)}=0.
\]
由定义
\[
    p_n=\dfrac{u^{(n)}}{2^n n!},
    \qquad
    p_n'=\dfrac{u^{(n+1)}}{2^n n!},
    \qquad
    p_n''=\dfrac{u^{(n+2)}}{2^n n!}.
\]
将上式除以 $2^n n!$ 并整体乘以 $-1$，即得
\[
    (1-x^2)p_n''-2xp_n'+n(n+1)p_n=0.
\]
\end{proof}

\methodsection{(E) 利用 Taylor 展开}

\begin{example}
求 $f(x)=\arctan x$ 在 $x=0$ 处的各阶导数.
\end{example}

\begin{proof}[解]
在 $|x|<1$ 内，
\[
    \arctan x
    =\sum\limits_{m=0}^{\infty}(-1)^m\dfrac{x^{2m+1}}{2m+1}.
\]
因而偶数阶导数在原点都为零，而
\[
    f^{(2m+1)}(0)
    =(2m+1)!\dfrac{(-1)^m}{2m+1}
    =(-1)^m(2m)!,
    \qquad m=0,1,2,\ldots.
\]
\end{proof}

\begin{example}
设
\[
f(x)=
\begin{cases}
    \dfrac{\sin x}{x}, & x\neq0,\\
    1, & x=0,
\end{cases}
\]
求 $f^{(k)}(0)$.
\end{example}

\begin{proof}[解]
在原点附近，
\[
    f(x)=\dfrac{\sin x}{x}
    =\sum\limits_{m=0}^{\infty}(-1)^m\dfrac{x^{2m}}{(2m+1)!}.
\]
故当 $k$ 为奇数时 $f^{(k)}(0)=0$；当 $k=2m$ 时，
\[
    f^{(2m)}(0)
    =(2m)!\dfrac{(-1)^m}{(2m+1)!}
    =\dfrac{(-1)^m}{2m+1}.
\]
即
\[
    f^{(k)}(0)=
    \begin{cases}
        0, & k\text{ 为奇数},\\
        \dfrac{(-1)^{k/2}}{k+1}, & k\text{ 为偶数}.
    \end{cases}
\]
\end{proof}

\subsection{利用导数求和}

\begin{example}
试求下列极限:
\begin{enumerate}[label=(\arabic*)]
    \item $\displaystyle
    I=\lim\limits_{x\to0}\dfrac{1}{x}\sum\limits_{i=1}^k f\left(\dfrac{x}{i}\right)$,
    已知 $f'(0)$ 存在,且 $f(0)=0,k\in\mathbb{N}$;
    \item $\displaystyle
    I=\lim\limits_{n\to\infty}n\left[\sum\limits_{i=1}^k f\left(a+\dfrac{i}{n}\right)-kf(a)\right]$,
    已知 $f'(a)$ 存在;
    \item $\displaystyle
    I=\lim\limits_{n\to\infty}
    \dfrac{\left(a+\dfrac{1}{n}\right)^n\left(a+\dfrac{2}{n}\right)^n\cdots\left(a+\dfrac{k}{n}\right)^n}{a^{nk}}$,
    $k\in\mathbb{N}$.
\end{enumerate}
\end{example}

\begin{proof}[解]
\begin{enumerate}[label=(\arabic*)]
    \item 由 $f(x)=f'(0)x+o(x)$，
    \[
        \dfrac{1}{x}\sum\limits_{i=1}^{k}f\left(\dfrac{x}{i}\right)
        =f'(0)\sum\limits_{i=1}^{k}\dfrac{1}{i}+o(1),
    \]
    故 $I=f'(0)\displaystyle\sum\limits_{i=1}^{k}1/i$。
    \item 对固定的 $i$，
    \[
        f\left(a+\dfrac{i}{n}\right)-f(a)
        =\dfrac{i}{n}f'(a)+o\left(\dfrac{1}{n}\right).
    \]
    对 $i=1,\ldots,k$ 求和得
    \[
        I=f'(a)\sum\limits_{i=1}^{k}i
        =\dfrac{k(k+1)}{2}f'(a).
    \]
    \item 假设 $a>0$。取对数后，
    \[
    \begin{aligned}
        \ln I_n
        &=n\sum\limits_{i=1}^{k}
        \ln\left(1+\dfrac{i}{an}\right)\\
        &\longrightarrow\dfrac{1}{a}\sum\limits_{i=1}^{k}i
        =\dfrac{k(k+1)}{2a}.
    \end{aligned}
    \]
    因而
    \[
        I=\exp\left(\dfrac{k(k+1)}{2a}\right).
    \]
\end{enumerate}
\end{proof}

\begin{example}
试求下列和式:
\begin{enumerate}[label=(\arabic*)]
    \item $\displaystyle I_n=\sum\limits_{k=0}^n ke^{kx}$;
    \item $\displaystyle \theta_n(x)=\sum\limits_{k=1}^n k^2x^{k-1}$;
    \item $\displaystyle I_n=\sum\limits_{k=1}^n C_{n}^{k}k^2$;
    \item $\displaystyle I_n=\sum\limits_{k=0}^{2n}(-1)^k C_{2n}^{k}k^n$;
    \item $\displaystyle S_n(x)=\sum\limits_{k=1}^n\dfrac{1}{2^k}\tan\left(\dfrac{x}{2^k}\right)$.
\end{enumerate}
\end{example}

\begin{proof}[解]
\begin{enumerate}[label=(\arabic*)]
    \item 令 $q=e^x$。对等比和求导，得
    \[
        I_n=\sum\limits_{k=0}^{n}kq^k
        =\dfrac{q-(n+1)q^{n+1}+nq^{n+2}}{(1-q)^2}.
    \]
    \item 由
    \[
        \sum\limits_{k=1}^{n}kx^k
        =\dfrac{x-(n+1)x^{n+1}+nx^{n+2}}{(1-x)^2},
    \]
    两端对 $x$ 求导得，当 $x\neq1$ 时
    \[
        \theta_n(x)=
        \dfrac{1+x-(n+1)^2x^n+(2n^2+2n-1)x^{n+1}-n^2x^{n+2}}
        {(1-x)^3}.
    \]
    当 $x=1$ 时，
    \[
        \theta_n(1)=\sum\limits_{k=1}^{n}k^2
        =\dfrac{n(n+1)(2n+1)}{6}.
    \]
    \item 对 $(1+x)^n$ 求导一次、两次，并在 $x=1$ 处取值，得
    \[
    \begin{aligned}
        \sum\limits_{k=1}^{n}C_k^n k^2
        &=\sum\limits_{k=1}^{n}C_k^n k(k-1)
        +\sum\limits_{k=1}^{n}C_k^n k\\
        &=n(n-1)2^{n-2}+n2^{n-1}\\
        &=n(n+1)2^{n-2}.
    \end{aligned}
    \]
    \item 该和是多项式 $P(k)=k^n$ 的 $2n$ 阶有限差（相差一个符号）。由于 $\deg P=n<2n$，其 $2n$ 阶有限差为零，故
    \[
        I_n=0.
    \]
    \item 使用恒等式
    \[
        \tan u=\cot u-2\cot2u.
    \]
    于是
    \[
    \begin{aligned}
        S_n(x)
        &=\sum\limits_{k=1}^{n}\left[
        \dfrac{1}{2^k}\cot\dfrac{x}{2^k}
        -\dfrac{1}{2^{k-1}}\cot\dfrac{x}{2^{k-1}}
        \right]\\
        &=\dfrac{1}{2^n}\cot\dfrac{x}{2^n}-\cot x.
    \end{aligned}
    \]
\end{enumerate}
\end{proof}

\begin{example}
设 $f_1(x),f_2(x),\ldots,f_n(x)$ 在 $x=x_0$ 处可导,且 $f_k(x_0)\neq0\ (k=1,2,\ldots,n)$,
\[
    F(x)=\prod\limits_{k=1}^n f_k(x),
\]
则
\[
    \dfrac{F'(x_0)}{F(x_0)}
    =\sum\limits_{k=1}^n\dfrac{f_k'(x_0)}{f_k(x_0)}.
\]
\end{example}

\begin{proof}[解]
由有限个函数乘积的求导法则，
\[
    F'(x_0)=\sum\limits_{j=1}^{n}f_j'(x_0)
    \prod\limits_{\substack{1\leq k\leq n\\k\neq j}}f_k(x_0).
\]
因为所有 $f_k(x_0)$ 都非零，两端除以
$F(x_0)=\displaystyle\prod\limits_{k=1}^{n}f_k(x_0)$，即得
\[
    \dfrac{F'(x_0)}{F(x_0)}
    =\sum\limits_{j=1}^{n}\dfrac{f_j'(x_0)}{f_j(x_0)}.
\]
\end{proof}

\subsection{利用导数证恒等式}

\begin{example}
证明:
\[
    2\arctan x+\arcsin\dfrac{2x}{1+x^2}=\pi,
    \qquad 1<x<\infty.
\]
\end{example}

\begin{proof}[解]
令
\[
    F(x)=2\arctan x+arcsin\dfrac{2x}{1+x^2}.
\]
当 $x>1$ 时，
\[
    \sqrt{1-\left(\dfrac{2x}{1+x^2}\right)^2}
    =\dfrac{x^2-1}{1+x^2}.
\]
因而
\[
    F'(x)=\dfrac{2}{1+x^2}
    +\dfrac{2(1-x^2)}{(1+x^2)^2}
    \dfrac{1+x^2}{x^2-1}=0.
\]
故 $F$ 在 $(1,+\infty)$ 上为常数。令 $x\to1^+$，得
\[
    F(x)\longrightarrow2\cdot\dfrac{\pi}{4}+\dfrac{\pi}{2}=\pi.
\]
因此题设恒等式成立。
\end{proof}

\begin{example}
设 $f(x),g(x)$ 在 $(-\infty,+\infty)$ 上均不是常数,可导且有
\[
\begin{cases}
    f(x+y)=f(x)f(y)-g(x)g(y),\\
    g(x+y)=f(x)g(y)+g(x)f(y),
\end{cases}
\]
其中 $x,y\in(-\infty,+\infty)$,若 $f'(0)=0$,则
\[
    f^2(x)+g^2(x)=1.
\]
\end{example}

\begin{proof}[解]
在加法公式中令 $y=0$。由于 $f,g$ 不同时为常数，可得
\[
    f(0)=1,
    \qquad
    g(0)=0.
\]
对两个加法公式关于 $y$ 求导，再令 $y=0$。记 $c=g'(0)$，并使用 $f'(0)=0$，得
\[
    f'(x)=-c g(x),
    \qquad
    g'(x)=c f(x).
\]
因此
\[
    \dfrac{\mathrm{d}}{\mathrm{d}x}\bigl[f(x)^2+g(x)^2\bigr]
    =2ff'+2gg'=0.
\]
该和为常数，并且在 $x=0$ 处等于 $1$，故
\[
    f(x)^2+g(x)^2=1.
\]
\end{proof}

\begin{example}
设 $f(x),g(x)$ 在 $(a,b)$ 上均不是常数,若有
\[
    f(x)+g(x)\neq0,
    \qquad
    f(x)g'(x)-f'(x)g(x)=0,
    \qquad x\in(a,b),
\]
则 $g(x)\neq0\ (x\in(a,b))$ 且
\[
    \dfrac{f(x)}{g(x)}\equiv C\text{(常数)}.
\]
\end{example}

\begin{proof}[解]
由 $f+g\neq0$，函数
\[
    H(x)=\dfrac{f(x)}{f(x)+g(x)}
\]
在 $(a,b)$ 上有定义。题设等式给出
\[
    H'(x)=\dfrac{f'(x)g(x)-f(x)g'(x)}{[f(x)+g(x)]^2}=0.
\]
故 $H(x)\equiv A$ 为常数，从而
\[
    \dfrac{g(x)}{f(x)+g(x)}\equiv1-A.
\]
若某点 $g(x_0)=0$，则 $1-A=0$，上式迫使 $g\equiv0$，与 $g$ 不是常数矛盾。因而 $g(x)\neq0$。又 $1-A\neq0$，所以
\[
    \dfrac{f(x)}{g(x)}
    =\dfrac{A}{1-A}
\]
为常数。
\end{proof}

\subsection{导函数的性质}

\begin{example}
设函数 $f(x)$ 在区间 $(-\infty,+\infty)$ 上二次可微,且有界,试证存在点 $x_0\in(-\infty,+\infty)$,使得 $f'(x_0)=0$.
\end{example}

\begin{proof}[解]
本题按当前题面不成立。取
\[
    f(x)=\arctan x.
\]
则 $f$ 在实轴上二次可微且有界，但
\[
    f'(x)=\dfrac{1}{1+x^2}>0
\]
对所有 $x\in\mathbb{R}$ 成立，因而不存在 $f'(x_0)=0$ 的点。原题需要补充条件或核对结论。
\end{proof}

\begin{example}
设函数 $f$ 在 $(0,1)$ 有界、可导,但 $\displaystyle\lim\limits_{x\to0^+}f(x)$ 不存在,试证存在数列 $x_n\to0^+\ (n\to\infty)$,使得
\[
    \lim\limits_{n\to+\infty}f'(x_n)=0.
\]
\end{example}

\begin{proof}[解]
反证。若不存在所求数列，则存在 $\varepsilon>0$ 与 $\delta>0$，使
\[
    |f'(x)|\geq\varepsilon,
    \qquad 0<x<\delta.
\]
导函数具有 Darboux 介值性，因而 $f'$ 在 $(0,\delta)$ 上不能变号。若 $f'\geq\varepsilon$，则对 $0<x<y<\delta$，中值定理给出
\[
    f(y)-f(x)\geq\varepsilon(y-x),
\]
故 $f$ 随 $x\to0^+$ 单调且有界，从而存在有限右极限；$f'\leq-\varepsilon$ 时同理。这与题设矛盾。

因此对每个 $n$，都能取 $0<x_n<1/n$ 使 $|f'(x_n)|<1/n$。于是 $x_n\to0^+$ 且 $f'(x_n)\to0$。
\end{proof}

\begin{example}
\begin{enumerate}[label=(\arabic*)]
    \item 设 $f(x)$ 在 $(a,b)$ 上可导,且有 $x_i\in(a,b),\lambda_i>0\ (i=1,2,\ldots,n)$,
    \[
        \sum\limits_{i=1}^n\lambda_i=1,
    \]
    则存在 $\xi\in(a,b)$,使得
    \[
        \sum\limits_{i=1}^n\lambda_i f'(x_i)=f'(\xi);
    \]
    \item 设 $f(x)$ 在 $(a,b)$ 上可导,且有 $x_i<y_i$,$x_i,y_i\in(a,b)\ (i=1,2,\ldots,n)$,则存在 $\xi\in(a,b)$,使得
    \[
        \sum\limits_{i=1}^n[f(y_i)-f(x_i)]
        =f'(\xi)\sum\limits_{i=1}^n(y_i-x_i).
    \]
\end{enumerate}
\end{example}

\begin{proof}[解]
\begin{enumerate}[label=(\arabic*)]
    \item 数 $\displaystyle\sum\limits_{i=1}^{n}\lambda_i f'(x_i)$ 是 $f'(x_i)$ 的凸组合，因而位于这些值的最小值与最大值之间。由导函数的 Darboux 介值性，存在 $\xi\in(a,b)$ 使
    \[
        f'(\xi)=\sum\limits_{i=1}^{n}\lambda_i f'(x_i).
    \]
    \item 由 Lagrange 中值定理，对每个 $i$ 存在 $c_i\in(x_i,y_i)$，使
    \[
        f(y_i)-f(x_i)=f'(c_i)(y_i-x_i).
    \]
    令
    \[
        \lambda_i=\dfrac{y_i-x_i}{\displaystyle\sum\limits_{j=1}^{n}(y_j-x_j)},
    \]
    则 $\lambda_i>0$ 且其和为 $1$。对 $f'(c_i)$ 应用第 (1) 问，即得题设等式。
\end{enumerate}
\end{proof}

\begin{example}
设 $f(x)$ 在 $\lbrack a,b\rbrack$ 上二次可导,若 $f''(x)\neq0\ (a\leq x\leq b)$,则对任意 $\xi\in(a,b)$,存在 $\xi'$ 满足 $a<\xi<\xi'\leq b$,使得
\[
    f'(\xi)=\dfrac{f(\xi')-f(a)}{\xi'-a}.
\]
\end{example}

\begin{proof}[解]
本题按当前题面不成立。取 $\lbrack a,b\rbrack=\lbrack0,1\rbrack$ 与
\[
    f(x)=e^x.
\]
则 $f''(x)=e^x\neq0$。若取 $\xi$ 充分接近 $1$，例如 $\xi=9/10$，则
\[
    f'(\xi)=e^{9/10}>e-1.
\]
但对任意 $\xi'\in(\xi,1\rbrack$，
\[
    \dfrac{f(\xi')-f(0)}{\xi'}
    =\dfrac{e^{\xi'}-1}{\xi'}
    \leq e-1<f'(\xi).
\]
因而不存在题设的 $\xi'$。原题中“对任意 $\xi$”或割线端点条件需要核对。
\end{proof}

\subsection{广义导数 *}

\begin{example}[\markstar\ Schwarz 定理]
若 $f(x)$ 在 $\lbrack a,b\rbrack$ 上连续,$f(x)$ 的广义二阶导数 $f^{[2]}(x)$ 存在且恒为零,则 $f(x)$ 是一次线性函数,即 $f(x)=ax+b$.
\end{example}

\begin{proof}[解]
构造辅助函数
\[
    F(x)=f(x)e^{g(x)}.
\]
由 $f(a)=f(b)=0$，有 $F(a)=F(b)=0$。根据 Rolle 定理，存在 $\xi\in(a,b)$ 使 $F'(\xi)=0$。又
\[
    F'(x)=e^{g(x)}\bigl[f'(x)+g'(x)f(x)\bigr].
\]
由于 $e^{g(\xi)}>0$，必有
\[
    f'(\xi)+g'(\xi)f(\xi)=0.
\]
\end{proof}

\begin{example}
若 $f(x)$ 在 $\lbrack a,b\rbrack$ 上连续,$f(a)<f(b)$,又设对一切 $x\in(a,b)$,
\[
    \lim\limits_{t\to0}\dfrac{f(x+t)-f(x-t)}{t}
\]
存在,用 $g(x)$ 表示这一极限值,则存在 $c\in(a,b)$,使得 $g(c)\geq0$.
\end{example}

\begin{proof}[解]
% TODO: 补充本题解答。
\end{proof}

\newpage

\section{微分中值定理}

\subsection{构造辅助函数(中值问题)}

\begin{example}
设 $f\in C(\lbrack a,b\rbrack)$,$g\in C(\lbrack a,b\rbrack)$,且均在 $(a,b)$ 上可导,若 $f(a)=f(b)=0$,则存在 $\xi\in(a,b)$,使得
\[
    f'(\xi)+g'(\xi)f(\xi)=0.
\]
\end{example}

\begin{proof}[解]
% TODO: 补充本题解答。
\end{proof}

\begin{example}
设 $f(x)$ 在 $\lbrack 0,1\rbrack$ 上二次可导,且 $f(0)=0$,则存在 $\xi\in(0,1)$,使得
\[
    f'(\xi)-(\xi-1)^2f''(\xi)=0.
\]
\end{example}

\begin{proof}[解]
% TODO: 补充本题解答。
\end{proof}

\begin{example}
设 $f(x),g(x)$ 在 $\lbrack a,b\rbrack$ 二次可导,且 $g''(x)\neq0$,若有
\[
    f(a)=f(b)=g(a)=g(b)=0,
\]
则存在 $\xi\in(a,b)$,使得
\[
    \dfrac{f(\xi)}{g(\xi)}=\dfrac{f''(\xi)}{g''(\xi)}.
\]
\end{example}

\begin{proof}[解]
% TODO: 补充本题解答。
\end{proof}

\begin{example}
设 $f(x),g(x)$ 在 $\lbrack a,b\rbrack$ 上可导,且在 $(a,b)$ 上 $g'(x)\neq0$,则存在 $\xi\in(a,b)$,使得
\[
    \dfrac{f(a)-f(\xi)}{g(\xi)-g(b)}=\dfrac{f'(\xi)}{g'(\xi)}.
\]
\end{example}

\begin{proof}[解]
% TODO: 补充本题解答。
\end{proof}

\begin{example}
设 $f\in C(\lbrack 0,1\rbrack)$,在 $(0,1)$ 上可导且 $f(1)=0$,则存在 $\xi\in(0,1)$ 使得
\[
    f'(\xi)=\left(1-\dfrac{1}{\xi}\right)f(\xi).
\]
\end{example}

\begin{proof}[解]
% TODO: 补充本题解答。
\end{proof}

\begin{example}
设 $f\in C(\lbrack a,b\rbrack)$,在 $(a,b)$ 上可导,若有
\[
    f(a)f(b)>0,
    \qquad
    f(a)f\left(\dfrac{a+b}{2}\right)<0,
\]
则对任给的 $\lambda\in(-\infty,+\infty)$,存在 $\xi\in(a,b)$,使得
\[
    f'(\xi)=\lambda f(\xi).
\]
\end{example}

\begin{proof}[解]
% TODO: 补充本题解答。
\end{proof}

\begin{example}
设 $f(x)$ 在 $\lbrack 0,+\infty)$ 上可导,若
\[
    f^2(a)-f^2(b)=b^2-a^2,
\]
则存在 $\xi\in(0,+\infty)$,使得
\[
    f'(\xi)=\dfrac{1-\xi^2}{(1+\xi^2)^2}.
\]
\end{example}

\begin{proof}[解]
% TODO: 补充本题解答。
\end{proof}

\begin{example}
设 $f(x)$ 在 $\lbrack 0,1\rbrack$ 上二次可导,且 $f(0)=f(1)=0$,则存在 $\xi\in(0,1)$,使得
\[
    f''(\xi)=\dfrac{2f'(\xi)}{1-\xi}.
\]
\end{example}

\begin{proof}[解]
% TODO: 补充本题解答。
\end{proof}

\begin{example}
设 $f(x)$ 在 $\lbrack 0,1\rbrack$ 上可导,且 $f(0)=0,f(x)>0\ (0<x<1)$,则存在 $\xi\in(0,1)$,使得
\[
    \dfrac{2f'(\xi)}{f(\xi)}
    =\dfrac{f'(1-\xi)}{f(1-\xi)}.
\]
\end{example}

\begin{proof}[解]
% TODO: 补充本题解答。
\end{proof}

\subsection{待定常数法(中值问题)}

\begin{example}
设 $f(x)$ 在 $\lbrack a,b\rbrack$ 上二次可导,则存在 $\xi\in(a,b)$,使得
\[
    f(b)-2f\left(\dfrac{a+b}{2}\right)+f(a)
    =\dfrac{(b-a)^2}{4}f''(\xi).
\]
\end{example}

\begin{proof}[解]
% TODO: 补充本题解答。
\end{proof}

\begin{example}
设 $f(x),g(x)$ 在 $\lbrack a,b\rbrack$ 上二次可导,则存在 $\xi\in(a,b)$,使得
\[
    \dfrac{f(b)-f(a)-(b-a)f'(a)}{g(b)-g(a)-(b-a)g'(a)}
    =\dfrac{f''(\xi)}{g''(\xi)}.
\]
\end{example}

\begin{proof}[解]
% TODO: 补充本题解答。
\end{proof}

\begin{example}
若存在 $f''(0)$,则
\[
    \lim\limits_{h\to0}
    \dfrac{f(2h)-2f(0)+f(-2h)}{4h^2}
    =f''(0).
\]
\end{example}

\begin{proof}[解]
% TODO: 补充本题解答。
\end{proof}

\begin{example}
设 $h>0$,$f\in C(\lbrack a-h,a+h\rbrack)$,在 $(a-h,a+h)$ 上可导,则存在 $\theta:0<\theta<1$,使得
\[
    \dfrac{f(a+h)-f(a-h)}{h}
    =f'(a+\theta h)+f'(a-\theta h).
\]
\end{example}

\begin{proof}[解]
% TODO: 补充本题解答。
\end{proof}

\begin{example}
设 $f(x)$ 在 $\lbrack a,b\rbrack$ 上连续,且在 $(a,b)$ 上二次可导,则对 $x\in(a,b)$,存在 $\xi\in(a,b)$,使得
\[
    \dfrac{f(x)-f(a)}{x-a}
    -\dfrac{f(b)-f(a)}{b-a}
    =\dfrac{1}{2}(x-b)f''(\xi).
\]
\end{example}

\begin{proof}[解]
% TODO: 补充本题解答。
\end{proof}

\begin{example}
设 $f(x)$ 在 $\lbrack a,c\rbrack$ 上二次可导,$a<b<c$,则存在 $\xi\in(a,c)$,使得
\[
    \dfrac{f(a)}{(a-b)(a-c)}
    +\dfrac{f(b)}{(b-c)(b-a)}
    +\dfrac{f(c)}{(c-a)(c-b)}
    =\dfrac{1}{2}f''(\xi).
\]
\end{example}

\begin{proof}[解]
% TODO: 补充本题解答。
\end{proof}

\begin{example}
设 $f(x)$ 在 $(a,b)$ 上三次可导,则存在 $\xi\in(a,b)$,使得
\[
    f(a)-f(b)+\dfrac{1}{2}(b-a)[f'(a)+f'(b)]
    =f'''(\xi)\dfrac{(b-a)^3}{12}.
\]
\end{example}

\begin{proof}[解]
% TODO: 补充本题解答。
\end{proof}

\subsection{利用几何意义(中值问题)}

\begin{example}
设 $f(x)$ 是可微函数,导函数 $f'(x)$ 严格单调递增,若 $f(a)=f(b)\ (a<b)$,试证明:对一切 $x\in(a,b)$,有
\[
    f(x)<f(a)=f(b).
\]
\end{example}

\begin{proof}[解]
% TODO: 补充本题解答。
\end{proof}

\begin{example}
设 $f(x)$ 在区间 $\lbrack 0,1\rbrack$ 上可微,$f(0)=0,f(1)=1$,$k_1,k_2,\ldots,k_n$ 为 $n$ 个正数,证明在区间 $\lbrack 0,1\rbrack$ 上存在一组互不相等的数 $x_1,x_2,\ldots,x_n$,使得
\[
    \sum\limits_{i=1}^n\dfrac{k_i}{f'(x_i)}
    =\sum\limits_{i=1}^n k_i.
\]
\end{example}

\begin{proof}[解]
% TODO: 补充本题解答。
\end{proof}

\begin{example}
设函数 $f(x)$ 在闭区间 $\lbrack a,b\rbrack$ 上连续,在开区间 $(a,b)$ 内可导,又 $f(x)$ 不是线性函数,且 $f(b)>f(a)$,试证:存在 $\xi\in(a,b)$,使得
\[
    f'(\xi)>\dfrac{f(b)-f(a)}{b-a}.
\]
\end{example}

\begin{proof}[解]
% TODO: 补充本题解答。
\end{proof}

\begin{example}
设 $f(x)$ 在 $\lbrack a,b\rbrack$ 上连续,在 $(a,b)$ 内可导,$f(x)$ 既不是常数也不是一次函数,试证:存在 $\xi\in(a,b)$,使得
\[
    |f'(\xi)|>\left|\dfrac{f(b)-f(a)}{b-a}\right|.
\]
\end{example}

\begin{proof}[解]
% TODO: 补充本题解答。
\end{proof}

\subsection{多中值问题}

\begin{example}
设 $f(x)$ 在 $\lbrack a,b\rbrack$ 上连续,在 $(a,b)$ 内可导,$f(a)\neq f(b)$,证明:存在 $\xi,\eta\in(a,b)$ 使得
\[
    f'(\xi)=\dfrac{a+b}{2\eta}f'(\eta).
\]
\end{example}

\begin{proof}[解]
% TODO: 补充本题解答。
\end{proof}

\begin{example}
设 $f(x)$ 在 $\lbrack a,b\rbrack\ (b>a>0)$ 上连续,在 $(a,b)$ 内可导,证明:在 $(a,b)$ 内存在两点 $\xi,\eta$,使得
\[
    f'(\xi)=\dfrac{\eta^2}{ab}f'(\eta).
\]
\end{example}

\begin{proof}[解]
% TODO: 补充本题解答。
\end{proof}

\begin{example}
设 $f(x)$ 在 $\lbrack a,b\rbrack$ 上连续,在 $(a,b)$ 内可导,$0<a<b$,证明:在 $(a,b)$ 内存在 $x_1,x_2,x_3$,使得
\[
    \dfrac{f'(x_1)}{2x_1}
    =(b^2+a^2)\dfrac{f'(x_2)}{4x_3^2}
    =\dfrac{\ln\dfrac{b}{a}}{b^2-a^2}x_3f'(x_3).
\]
\end{example}

\begin{proof}[解]
% TODO: 补充本题解答。
\end{proof}

\begin{example}
设 $f(x)$ 在区间 $\lbrack a,b\rbrack$ 上连续,在 $(a,b)$ 内可导,且 $a\geq0$(或 $b\leq0$),试证:
\begin{enumerate}[label=(\arabic*)]
    \item 存在 $x_1,x_2,x_3\in(a,b)$,使得
    \[
        f'(x_1)
        =(b+a)\dfrac{f'(x_2)}{2x_2}
        =(b^2+ba+a^2)\dfrac{f'(x_3)}{3x_3^2};
    \]
    \item 对任意 $n\in\{1,2,\ldots\}$,存在 $x_1,x_2,\ldots,x_n\in(a,b)$,使得
    \[
    \begin{aligned}
        f'(x_1)
        &=(b+a)\dfrac{f'(x_2)}{2x_2}
        =\cdots\\
        &=(b^{n-1}+b^{n-2}a+\cdots+ba^{n-2}+a^{n-1})
        \dfrac{f'(x_n)}{nx_n^{n-1}}.
    \end{aligned}
    \]
\end{enumerate}
\end{example}

\begin{proof}[解]
% TODO: 补充本题解答。
\end{proof}

\begin{example}
设 $f\in C(\lbrack a,b\rbrack)$,且在 $(a,b)$ 上可导,若 $f(a)=f(b)=1$,则存在 $\xi,\eta\in(a,b)$,使得
\[
    e^{\eta-\xi}[f(\eta)+f'(\eta)]=1.
\]
\end{example}

\begin{proof}[解]
% TODO: 补充本题解答。
\end{proof}

\begin{example}
设 $f\in C(\lbrack a,b\rbrack)$,且在 $(a,b)$ 上可导,若 $f(0)=0,f(1)=1$,则存在 $\xi,\eta\in(a,b)$,使得
\[
    f'(\xi)f'(\eta)=1,
    \qquad \xi\neq\eta.
\]
\end{example}

\begin{proof}[解]
% TODO: 补充本题解答。
\end{proof}

\subsection{函数与导函数的关系}

\begin{example}
设 $f(x)$ 在 $(a,+\infty)$ 上可导,若 $\displaystyle\lim\limits_{x\to+\infty}f'(x)=B>0$,则:
\begin{enumerate}[label=(\arabic*)]
    \item $\displaystyle\lim\limits_{x\to+\infty}f(x)=+\infty$;
    \item $f(x)\sim Bx\ (x\to+\infty)$.
\end{enumerate}
\end{example}

\begin{proof}[解]
% TODO: 补充本题解答。
\end{proof}

\begin{example}[\markstar]
设 $f\in C^{(1)}((a,+\infty))$,若 $f'(x)$ 在 $(a,+\infty)$ 上一致连续,且存在 $\displaystyle\lim\limits_{x\to+\infty}f(x)$,则
\[
    \lim\limits_{x\to+\infty}f'(x)=0.
\]
\end{example}

\begin{proof}[解]
% TODO: 补充本题解答。
\end{proof}

\begin{example}[\markstar]
设 $f(x)$ 在 $\lbrack 0,+\infty)$ 上三阶可导,且有
\[
    \lim\limits_{x\to+\infty}f(x)=A,
    \qquad
    \lim\limits_{x\to+\infty}f'''(x)=0,
\]
则
\[
    \lim\limits_{x\to+\infty}f'(x)
    =\lim\limits_{x\to+\infty}f''(x)=0.
\]
\end{example}

\begin{proof}[解]
% TODO: 补充本题解答。
\end{proof}

\begin{example}[\markstar]
试证明下列命题:
\begin{enumerate}[label=(\arabic*)]
    \item 设 $f(x)$ 在 $(0,+\infty)$ 上可导,$a>0$,若
    \[
        \lim\limits_{x\to+\infty}[af(x)+f'(x)]=l,
    \]
    则
    \[
        \lim\limits_{x\to+\infty}f(x)=\dfrac{l}{a};
    \]
    \item 设 $f(x)$ 在 $(0,+\infty)$ 上可导,$a>0$,若有
    \[
        \lim\limits_{x\to+\infty}[af(x)+2\sqrt{x}f'(x)]=l,
    \]
    则
    \[
        \lim\limits_{x\to+\infty}f(x)=\dfrac{l}{a};
    \]
    \item 设 $f(x)$ 在 $(0,+\infty)$ 上二次可导,若有
    \[
        \lim\limits_{x\to+\infty}[f(x)+f'(x)+f''(x)]=l,
    \]
    则
    \[
        \lim\limits_{x\to+\infty}f(x)=l.
    \]
\end{enumerate}
\end{example}

\begin{proof}[解]
% TODO: 补充本题解答。
\end{proof}

\begin{example}
设 $f(x)$ 在 $(0,+\infty)$ 上可导,若 $\displaystyle\lim\limits_{x\to+\infty}f'(x)=l$,则对任意的 $A$ 值,有
\[
    \lim\limits_{x\to+\infty}[f(x+A)-f(x)]=lA.
\]
\end{example}

\begin{proof}[解]
% TODO: 补充本题解答。
\end{proof}

\begin{example}[\markstar]
设 $f(x)$ 在 $\lbrack 0,+\infty)$ 上可微,$f(0)=0$,且
\[
    |f'(x)|\leq A|f(x)|,
    \qquad x\in\lbrack 0,+\infty),
\]
则 $f\equiv0$.
\end{example}

\begin{proof}[解]
% TODO: 补充本题解答。
\end{proof}

\begin{example}
设 $f(x),g(x)$ 在 $\lbrack a,b\rbrack$ 上连续,$g(x)$ 在 $(a,b)$ 内可微,$g(a)=0$,若有实数 $\lambda\neq0$,使得
\[
    |g(x)f(x)+\lambda g'(x)|\leq|g(x)|,
    \qquad x\in(a,b),
\]
成立,则 $g\equiv0$.
\end{example}

\begin{proof}[解]
% TODO: 补充本题解答。
\end{proof}

\begin{example}
证明:在 $\mathbb{R}$ 上的任何可微函数都不可能满足下列函数方程:
\begin{enumerate}[label=(\arabic*)]
    \item $g(g(x))=-x^3+x+1$;
    \item $g(g(x))=-x^3+x^2\pm1$;
    \item $g(g(x))=x^2-3x+3$.
\end{enumerate}
\end{example}

\begin{proof}[解]
% TODO: 补充本题解答。
\end{proof}

\newpage

\section{Taylor 公式}

\subsection{中值等式}

\begin{example}
设 $f(x)$ 在 $\lbrack a,b\rbrack$ 上三次可导,试证:存在 $c\in(a,b)$,使得
\[
    f(b)=f(a)+f'\left(\dfrac{a+b}{2}\right)(b-a)
    +\dfrac{1}{24}f'''(c)(b-a)^3.
\]
\end{example}

\begin{proof}[解]
% TODO: 补充本题解答。
\end{proof}

\begin{example}
设 $f\in C^{(3)}(\lbrack -1,1\rbrack)$,且 $f(-1)=0,f(1)=1,f'(0)=0$,则存在 $\xi\in(-1,1)$,使得
\[
    f'''(\xi)=3.
\]
\end{example}

\begin{proof}[解]
% TODO: 补充本题解答。
\end{proof}

\begin{example}
设 $f(x)$ 在 $\lbrack -1,1\rbrack$ 上二次可导,若有 $f(-1)=0,f(0)=0,f(1)=1$,则存在 $\xi\in(-1,1)$,使得
\[
    f''(\xi)=1.
\]
\end{example}

\begin{proof}[解]
% TODO: 补充本题解答。
\end{proof}

\begin{example}
设 $f(x)$ 在 $\lbrack -1,1\rbrack$ 上三次可导,若有 $f(-1)=f(0)=0,f(1)=1,f'(0)=0$,则存在 $\xi\in(-1,1)$,使得
\[
    f'''(\xi)\geq 3.
\]
\end{example}

\begin{proof}[解]
% TODO: 补充本题解答。
\end{proof}

\subsection{界的估计}

\begin{example}
设 $f(x)$ 二次可微,且 $f(0)=f(1)=0$,又
$\displaystyle\max\limits_{0\leq x\leq1}f(x)=2$,试证:
\[
    \min\limits_{0\leq x\leq1}f''(x)\leq-16.
\]
\end{example}

\begin{proof}[解]
% TODO: 补充本题解答。
\end{proof}

\begin{example}
已知 $f(x)$ 在 $\lbrack 0,1\rbrack$ 上二阶可导,且 $f(0)=0,f(1)=3$,又
$\displaystyle\min\limits_{x\in\lbrack 0,1\rbrack}f(x)=-1$,证明:存在 $c\in(0,1)$,使得
\[
    f''(c)\geq18.
\]
\end{example}

\begin{proof}[解]
% TODO: 补充本题解答。
\end{proof}

\begin{example}[\markstar]
设 $f(x)$ 在 $\lbrack 0,1\rbrack$ 上二次可导,若有
\[
    |f(x)|\leq A,\qquad |f''(x)|\leq B,\qquad x\in\lbrack 0,1\rbrack,
\]
则
\[
    |f'(x)|\leq2A+\dfrac{B}{2}.
\]
\end{example}

\begin{proof}[解]
% TODO: 补充本题解答。
\end{proof}

\begin{example}
设 $f(x)$ 在 $\lbrack 0,1\rbrack$ 上二阶可导,且有 $f(0)=f(1)$,$|f''(x)|\leq M$,则
\[
    |f'(x)|\leq\dfrac{M}{2}.
\]
\end{example}

\begin{proof}[解]
% TODO: 补充本题解答。
\end{proof}

\begin{example}[\markstar]
设 $f(x)$ 在 $(-\infty,+\infty)$ 上二次可导,且有
\[
    M_k=\sup\{|f^{(k)}(x)|:-\infty<x<+\infty\}<+\infty.
\]
则
\[
    M_1\leq\sqrt{2M_0M_2}.
\]
\end{example}

\begin{proof}[解]
% TODO: 补充本题解答。
\end{proof}

\begin{example}[\markstar]
设 $f(x)$ 在 $(-\infty,+\infty)$ 上 $m$ 次可导,且有
\[
    M_k=\sup\{|f^{(k)}(x)|:-\infty<x<+\infty\}<+\infty.
\]
则
\[
    M_k\leq2^{\frac{k(m-k)}{2}}M_0^{1-\frac{k}{m}}M_m^{\frac{k}{m}},
    \qquad k=1,2,\ldots,m-1.
\]
\end{example}

\begin{proof}[解]
% TODO: 补充本题解答。
\end{proof}

\subsection{中值点的极限}

\begin{example}
证明:对每个 $x>0$,均存在 $\theta(x)\in\lbrack \dfrac{1}{4},\dfrac{1}{2}\rbrack$,使得
\[
    \sqrt{x+1}-\sqrt{x}=\dfrac{1}{2\sqrt{x+\theta(x)}}.
\]
且
\[
    \lim\limits_{x\to0^+}\theta(x)=\dfrac{1}{4},
    \qquad
    \lim\limits_{x\to+\infty}\theta(x)=\dfrac{1}{2}.
\]
\end{example}

\begin{proof}[解]
% TODO: 补充本题解答。
\end{proof}

\begin{example}
设 $f(x)$ 在点 $a$ 的某个邻域上 $n+2$ 阶连续可微,且 $f^{(n+2)}(a)\neq0$,并有
\[
    f(a+h)=f(a)+f'(a)h+\cdots+\dfrac{f^{(n)}(a)}{n!}h^n
    +\dfrac{f^{(n+1)}(a+\theta h)}{(n+1)!}h^{n+1},
\]
其中 $h\in(0,\delta)$,$0<\theta<1$.证明:
\[
    \lim\limits_{h\to0}\theta=\dfrac{1}{n+2}.
\]
\end{example}

\begin{proof}[解]
% TODO: 补充本题解答。
\end{proof}

\begin{example}
设 $f(x)$ 在 $(x_0-\delta,x_0+\delta)$ 上 $n$ 阶连续可微,已知
\[
    f^{(2)}(x_0)=f^{(3)}(x_0)=\cdots=f^{(n-1)}(x_0)=0,
    \qquad f^{(n)}(x_0)\neq0,
\]
若 $0<|h|<\delta$ 时有
\[
    \dfrac{f(x_0+h)-f(x_0)}{h}=f'(x_0+\theta h),
\]
证明:
\[
    \lim\limits_{h\to0}\theta
    =\sqrt[n-1]{\dfrac{1}{n}}.
\]
\end{example}

\begin{proof}[解]
% TODO: 补充本题解答。
\end{proof}

\begin{example}
设 $f(x)$ 在点 $a$ 处二阶可导,且 $f''(a)=0$.证明当 $h$ 充分小时,有等式
\[
    f(a+h)-f(a)=f'(a+\theta h)h,
\]
其中 $\theta\in(0,1)$,并且
\[
    \lim\limits_{h\to0}\theta=\dfrac{1}{2}.
\]
\end{example}

\begin{proof}[解]
% TODO: 补充本题解答。
\end{proof}

\subsection{极值问题}

\begin{example}
设 $f\in C(\lbrack a,b\rbrack)$,且 $f\in C^{(2)}((a,b))$,若有
\[
    f''(x)=e^x f(x),\qquad f(a)=f(b)=0,
\]
则 $f(x)\equiv0$.
\end{example}

\begin{proof}[解]
% TODO: 补充本题解答。
\end{proof}

\begin{example}
设 $f\in C^{(2)}((-\infty,+\infty))$,若有
\[
    [f(0)]^2+[f'(0)]^2=4,
    \qquad |f(x)|=1\quad(-\infty<x<+\infty),
\]
则存在 $\xi$,使得
\[
    f(\xi)+f''(\xi)=0.
\]
\end{example}

\begin{proof}[解]
% TODO: 补充本题解答。
\end{proof}

\begin{example}
设 $f(x),g(x)$ 定义在 $(-\infty,+\infty)$ 上,且 $f(x)$ 二次可导,又有
\[
    f(a)=f(b)=0\quad(a<b),
\]
且
\[
    f''(x)+f'(x)g(x)-f(x)=0,
    \qquad x\in(-\infty,+\infty),
\]
则 $f(x)\equiv0$.
\end{example}

\begin{proof}[解]
% TODO: 补充本题解答。
\end{proof}

\begin{example}
设 $f\in C^{(2)}(\lbrack 0,a\rbrack)$,且 $x_0\in(0,a)$ 是 $f(x)$ 的极小值点,则
\[
    |f'(0)|+|f'(a)|
    \leq a\max\limits_{x\in\lbrack 0,a\rbrack}|f''(x)|.
\]
\end{example}

\begin{proof}[解]
% TODO: 补充本题解答。
\end{proof}

\newpage

\section{凹凸性与不等式}

\subsection{凹凸性}

\begin{example}
设 $f(x)$ 为区间 $(a,b)$ 内的凸函数,试证:$f(x)$ 在 $I$ 的任一闭区间 $\lbrack \alpha,\beta\rbrack\subseteq(a,b)$ 上满足 Lipschitz 条件.
\end{example}

\begin{proof}[解]
% TODO: 补充本题解答。
\end{proof}

\begin{example}
设 $f(0)=0$,$f(x)$ 在 $\lbrack 0,+\infty)$ 上为非负的严格凸函数,令
\[
    F(x)=\dfrac{f(x)}{x},\qquad x>0,
\]
证明:$f(x),F(x)$ 为严格递增的.
\end{example}

\begin{proof}[解]
% TODO: 补充本题解答。
\end{proof}

\begin{example}
设 $f(x),g(x)$ 在区间 $(a,b)$ 内有定义,且对任意 $x,x_0\in(a,b)$ 有
\[
    f(x)-f(x_0)\geq g(x_0)(x-x_0).
\]
试证:对任意 $x_0\in(a,b)$,$f(x)$ 在 $x_0$ 处连续,并且有左右单侧导数.
\end{example}

\begin{proof}[解]
% TODO: 补充本题解答。
\end{proof}

\begin{example}
证明下列命题成立:
\begin{enumerate}[label=(\arabic*)]
    \item 两凸函数之和仍为凸函数;
    \item 两递增非负凸函数之积仍为凸函数.
\end{enumerate}
\end{example}

\begin{proof}[解]
% TODO: 补充本题解答。
\end{proof}

\begin{example}
设 $f\in C(\lbrack 0,1\rbrack)$,且在 $(0,1)$ 上二次可导,若有
\[
    f(0)=f(1)=0,
    \qquad f''(x)+2f'(x)+f(x)=0\quad(0<x<1),
\]
则 $f(x)\leq0$.
\end{example}

\begin{proof}[解]
% TODO: 补充本题解答。
\end{proof}

\begin{example}
设 $f(x)$ 在 $(a,b)$ 上四次可导,$x_0\in(a,b)$,且有
\[
    f^{(2)}(x_0)=0,
    \qquad f^{(3)}(x_0)=0,
    \qquad f^{(4)}(x_0)>0,
\]
则 $f(x)$ 是下凸函数.
\end{example}

\begin{proof}[解]
% TODO: 补充本题解答。
\end{proof}

\begin{example}
设 $f(x)$ 在 $(-\infty,+\infty)$ 上二次可导,且有
\[
    f(x)\leq0,
    \qquad f''(x)\geq0,
\]
则 $f(x)\equiv C$(常数).
\end{example}

\begin{proof}[解]
% TODO: 补充本题解答。
\end{proof}

\begin{example}
设 $f(x)$ 在 $(a,b)$ 上可导,若对 $(a,b)$ 中的 $x,y$($x\neq y$),存在唯一的 $\xi\in(a,b)$,使得
\[
    \dfrac{f(y)-f(x)}{y-x}=f'(\xi),
\]
则 $f(x)$ 是严格上凸或下凸的.
\end{example}

\begin{proof}[解]
% TODO: 补充本题解答。
\end{proof}

\begin{example}
设 $f(x)$ 是 $(0,1)$ 上三次可导的非负函数,若存在 $x_1,x_2\in(0,1)$ 使得 $f(x_1)=f(x_2)=0$,则存在点 $\xi\in(0,1)$,使得
\[
    f'''(\xi)=0.
\]
\end{example}

\begin{proof}[解]
% TODO: 补充本题解答。
\end{proof}

\subsection{不等式}

\methodsection{(A) 利用单调性}

\begin{example}
对任意 $n=1,2,3,\ldots$,求能使不等式
\[
    \left(1+\dfrac{1}{n}\right)^{n+\alpha}
    \leq e
    \leq\left(1+\dfrac{1}{n}\right)^{n+\beta}
\]
成立的 $\alpha$ 的最大值与 $\beta$ 的最小值.
\end{example}

\begin{proof}[解]
% TODO: 补充本题解答。
\end{proof}

\begin{example}
证明如下不等式:
\begin{enumerate}[label=(\arabic*)]
    \item 当 $0<x_1<x_2<\dfrac{\pi}{2}$ 时,有
    \[
        \dfrac{\tan x_2}{\tan x_1}>\dfrac{x_2}{x_1};
    \]
    \item 当 $x\in\left(0,\dfrac{\pi}{2}\right)$ 时,有
    \[
        \dfrac{\tan x}{x}>\dfrac{x}{\sin x};
    \]
    \item 当 $x\in\left(0,\dfrac{\pi}{2}\right)$ 时,有
    \[
        2\sin x+3\tan x>3x.
    \]
\end{enumerate}
\end{example}

\begin{proof}[解]
% TODO: 补充本题解答。
\end{proof}

\begin{example}
\begin{enumerate}[label=(\arabic*)]
    \item 设 $b>a>e$,证明 $a^b>b^a$;
    \item 比较 $\pi^e$ 与 $e^\pi$ 的大小.
\end{enumerate}
\end{example}

\begin{proof}[解]
% TODO: 补充本题解答。
\end{proof}

\methodsection{(B) 利用中值定理}

\begin{example}
证明:
\begin{enumerate}[label=(\arabic*)]
    \item 当 $x>0$ 时,
    \[
        \ln\left(1+\dfrac{1}{x}\right)>\dfrac{1}{x+1};
    \]
    \item 当 $x>0$ 时,
    \[
        \dfrac{x}{1+x^2}<\arctan x<x.
    \]
\end{enumerate}
\end{example}

\begin{proof}[解]
% TODO: 补充本题解答。
\end{proof}

\begin{example}
\begin{enumerate}[label=(\arabic*)]
    \item 证明:当 $s>0$ 时,
    \[
        \dfrac{n^{s+1}}{s+1}
        <1^s+2^s+\cdots+n^s
        <\dfrac{(n+1)^{s+1}}{s+1};
    \]
    \item 设 $\displaystyle\lambda=\sum\limits_{k=1}^n\dfrac{1}{k}$,证明
    \[
        e^\lambda>n+1.
    \]
\end{enumerate}
\end{example}

\begin{proof}[解]
% TODO: 补充本题解答。
\end{proof}

\methodsection{(C) 利用凹凸性}

\begin{example}
证明下列不等式:
\begin{enumerate}[label=(\arabic*)]
    \item 当 $x\in\lbrack 0,1\rbrack$ 时,
    \[
        1+x^2\leq2^x\leq1+x;
    \]
    \item 当 $x\in\lbrack 0,1\rbrack$ 时,
    \[
        \sin(\pi x)\leq\dfrac{\pi^2}{2}x(1-x).
    \]
\end{enumerate}
\end{example}

\begin{proof}[解]
% TODO: 补充本题解答。
\end{proof}

\begin{example}[\markstar]
证明下列经典不等式:
\begin{enumerate}[label=(\arabic*)]
    \item (平均值不等式)设 $a_i>0$,则
    \[
        \dfrac{n}{\displaystyle\sum\limits_{i=1}^n\dfrac{1}{a_i}}
        \leq\sqrt[n]{a_1a_2\cdots a_n}
        \leq\dfrac{\displaystyle\sum\limits_{i=1}^n a_i}{n}.
    \]
    \item (Young 不等式)设 $a,b>0$,$0<p,q<1$,且 $\dfrac{1}{p}+\dfrac{1}{q}=1$,则
    \[
        ab\leq\dfrac{a^p}{p}+\dfrac{b^q}{q}.
    \]
    \item (Holder 不等式)设 $a_i,b_i>0$,$0<p,q<1$,且 $\dfrac{1}{p}+\dfrac{1}{q}=1$,则
    \[
        \sum\limits_{i=1}^n a_i b_i
        \leq
        \left(\sum\limits_{i=1}^n a_i^p\right)^{\frac{1}{p}}
        \left(\sum\limits_{i=1}^n b_i^q\right)^{\frac{1}{q}}.
    \]
\end{enumerate}
\end{example}

\begin{proof}[解]
% TODO: 补充本题解答。
\end{proof}

\end{document}

